\scanpage{13}
\uline{Questions sur courbes elliptiques} sur $\overline Q$.

Soit $C$ une telle courbe, $G$ son groupe de Galois motivique [$G=\mathrm{Gl}(2,Q)$ ou $G=\prod_{K/Q}G_m$, $K$ le corps de multiplication complexe]. Soit $g\in G(C)/G(Q)$ définissant l'application du \unsure{réseau|viseur} transcendant de Betti dans celui de De Rham.

a) $Q(g)$ de degré de transcendance égal à $\dim G$ ($4$ ou $2$)?

b) Soit $T$ le lieu de $g$ sur $G_{\overline Q}$, \deleted{\unsure{$H\subset G$}} qu'on peut supposer (par translation à droite) contenir $e$, et soit $H$ le groupe engendré \deleted{et} par $T$. \deleted{et ses conjugués}
A-t-on $H=G$
\deleted{\unsure{algébriquement}}

c) Soit \deleted{$g\in H_{\overline Q}\subset G'_{\overline Q}$ \unsure{le sous-groupe connexe}} \inserted{$h\in C/\overline Q$} [groupe additif) définissant le splitting transcendant (via th. de Hodge) de De Rham. Peut-on avoir $h\in\overline Q$, ou $h$ est-il tjs transcendant sur $Q$ [\uline{cas où} $C$ \uline{n'a pas de multiplication complexe}]

\margin{\uline{NB} si $C$ est à multiplication complexe, on a tjs $h\in\overline Q$!}
\mspage{12}
